\chapter{绪论}\label{Chap:chap1}

\section{研究背景与动机}

随着云计算、大数据、人工智能和物联网技术的蓬勃发展，全球数据中心网络的规模和复杂度呈现指数级增长态势。根据国际数据公司（International Data Corporation，IDC）的统计数据，全球数据中心网络流量正以年均25\%以上的速度持续增长，预计到2026年将突破25ZB（泽字节）的惊人规模\upcite{idc2023}。传输控制协议（Transmission Control Protocol，TCP）作为互联网传输层的核心协议，承载了超过90\%的数据中心网络流量，其拥塞控制机制的性能直接影响着网络资源利用效率、应用响应延迟和用户服务质量。网络拥塞控制是计算机网络领域最具挑战性的研究课题之一，自互联网诞生以来一直是学术界和工业界持续关注的焦点问题。

在现代数据中心网络架构中，网络带宽经历了从千兆级别到10Gbps、25Gbps乃至100Gbps级别的跨越式提升，而网络延迟则从毫秒级降低到微秒级。以典型的Leaf-Spine数据中心网络架构为例，机架内（Intra-Rack）通信延迟通常在微秒级别，而跨Pod（Inter-Pod）通信延迟也仅在数十微秒量级。这种高带宽、低延迟的网络环境对传统拥塞控制算法提出了前所未有的严峻挑战。

现代数据中心承载着多种类型的应用工作负载，这些应用对网络性能有着截然不同的需求。分布式存储系统（如HDFS、Ceph）需要高吞吐量以支持大规模数据读写；在线事务处理系统（OLTP）和分布式数据库（如Spanner、CockroachDB）对延迟极为敏感，通常要求99百分位延迟（P99 Latency）控制在毫秒级甚至亚毫秒级；机器学习训练任务（如分布式TensorFlow、PyTorch）则对带宽利用率和多流公平性有较高要求。单一的拥塞控制策略难以同时满足这些多样化的性能需求。

更为复杂的是，数据中心网络的流量模式具有高度动态性和突发性特征。研究表明\upcite{datacentertraffic}，数据中心网络流量呈现出显著的"Mice-Elephant"（小流量与大流量并存）分布特征：绝大多数流量是持续时间短、数据量小的"Mice Flow"（小鼠流），而少数持续时间长、数据量大的"Elephant Flow"（大象流）则贡献了绝大部分的网络带宽消耗。Incast（多对一）通信模式在分布式系统中十分常见，会导致瞬时的严重拥塞；而Partition/Aggregate（分区聚合）工作负载则对尾延迟（Tail Latency）极为敏感。

传统的TCP拥塞控制算法在这种复杂的网络环境中显露出诸多不足。基于丢包的算法（如TCP NewReno\upcite{newreno}、TCP CUBIC\upcite{cubic}）将数据包丢失作为拥塞信号，但丢包是拥塞的滞后指示器，通常意味着网络缓冲区已经溢出，排队延迟已经积累到较高水平。基于延迟的算法（如TCP Vegas\upcite{vegas}、TCP BBR\upcite{bbr}）尝试通过监测RTT变化提前感知拥塞，但RTT测量易受多种因素干扰，且在与基于丢包算法共存时存在公平性问题。

显式拥塞通知（Explicit Congestion Notification，ECN）机制为解决上述问题提供了新的思路。ECN允许网络设备在队列长度超过阈值时标记数据包而非丢弃，从而实现更早、更精确的拥塞信号传递。数据中心TCP（DCTCP\upcite{dctcp}）利用ECN标记比例进行细粒度调整，取得了良好效果。然而，现有算法对ECN信号的利用仍不够充分，大多采用静态参数配置，缺乏对动态网络环境的自适应能力。

近年来，机器学习特别是强化学习技术的快速发展为网络拥塞控制带来了革命性的新思路\upcite{rlcc}。强化学习（Reinforcement Learning，RL）是一种通过与环境交互来学习最优策略的机器学习方法。在强化学习框架中，智能体（Agent）通过观察环境状态、执行动作并获得奖励来不断优化决策策略。将强化学习应用于网络拥塞控制，理论上能够实现：（1）自适应地学习网络环境特征，无需人工设计复杂的控制规则；（2）综合利用多种拥塞信号进行决策；（3）根据网络条件动态调整控制参数；（4）良好的可扩展性和通用性。

然而，现有基于强化学习的拥塞控制方案仍存在诸多问题：状态空间设计不完善，通常只使用4-6维的有限状态参数；奖励函数设计不合理，难以平衡多个优化目标；缺乏对ECN信号的分级利用；算法验证效率低下等。这些问题限制了强化学习在拥塞控制领域的实际应用效果。

\section{研究目标与主要贡献}

针对上述问题，本文提出了Lark算法——一种基于强化学习的多信号融合自适应网络拥塞控制方法。"Lark"意为"云雀"，云雀以其轻盈敏捷的飞行姿态和清脆悦耳的歌声著称。选用"Lark"作为算法名称，寓意本算法追求的设计目标：如同云雀般轻盈敏捷，实现高吞吐量（High Throughput）与低延迟（Low Latency）的完美平衡；如同云雀对环境变化的敏锐感知，实现对多种拥塞信号的快速响应和智能融合。

图\ref{fig:lark_contribution_overview}展示了Lark算法的核心贡献概览及各组件之间的逻辑关系。

\begin{figure}[htbp]
\centering
\begin{tikzpicture}[
    node distance=1.2cm and 1.5cm,
    block/.style={rectangle, draw, fill=blue!8, text width=5.8cm, minimum height=1.0cm, align=center, rounded corners=3pt, font=\small},
    arrow/.style={->, >=stealth, thick, draw=gray!70},
]
% 顶层目标
\node[block, fill=red!12, text width=10cm, minimum height=0.9cm] (goal) {Lark设计目标：高吞吐量（High Throughput） $+$ 低延迟（Low Latency）};

% 第二层：核心创新
\node[block, below left=1.0cm and 0.5cm of goal] (obs) {贡献1：15维观测空间\\{\scriptsize 首次纳入ECN状态、CA状态、拥塞事件}};
\node[block, below right=1.0cm and 0.5cm of goal] (detect) {贡献2：多信号融合拥塞检测\\{\scriptsize 分层优先级 $+$ 频率过滤}};

% 第三层
\node[block, below=1.0cm of obs] (response) {贡献3：差异化拥塞响应\\{\scriptsize 丢包0.70 / ECN 0.85 / 超时0.50}};
\node[block, below=1.0cm of detect] (adapt) {贡献4：自适应参数调优\\{\scriptsize $\alpha \in [1.05, 1.50]$，奖励EMA驱动}};

% 第四层
\node[block, below right=1.0cm and -2.2cm of response, text width=10cm] (fusion) {贡献5：融合控制策略 $V_t = \max(\alpha \times BDP, W) + \gamma \times MSS$};

% 底层
\node[block, below=0.8cm of fusion, fill=green!8, text width=10cm] (impl) {实现验证：ns-3 $+$ ns3-gym/OpenGym $+$ ZeroMQ IPC $+$ MTP并行仿真};

% 箭头
\draw[arrow] (goal.south) -- ++(0,-0.3) -| (obs.north);
\draw[arrow] (goal.south) -- ++(0,-0.3) -| (detect.north);
\draw[arrow] (obs) -- (response);
\draw[arrow] (detect) -- (adapt);
\draw[arrow] (response.south) -- ++(0,-0.3) -| (fusion.north);
\draw[arrow] (adapt.south) -- ++(0,-0.3) -| (fusion.north);
\draw[arrow] (fusion) -- (impl);
\end{tikzpicture}
\caption{Lark算法核心贡献概览}
\label{fig:lark_contribution_overview}
\end{figure}

本文的主要研究贡献如下：

\textbf{（1）完整的15维观测空间设计}

本文设计了一个包含15个参数的完整观测空间，首次将显式拥塞通知（ECN）状态、拥塞控制状态机状态、拥塞事件类型等丰富的TCP协议栈信息纳入强化学习状态表示。与现有方案通常采用的4-6维简化状态空间相比，Lark的15维观测空间能够为强化学习智能体提供更全面、更精细的网络状态感知能力，为精准的拥塞控制决策奠定基础。

观测空间的15个参数涵盖六大类别：连接标识信息（套接字UUID、节点ID）用于多流识别和状态隔离；时间信息（仿真时间戳）用于时序分析和事件率计算；窗口状态参数（慢启动阈值ssthresh、拥塞窗口cwnd）是拥塞控制的核心状态变量；传输性能指标（段大小、确认段数、在途字节数）用于吞吐量估算和管道利用率计算；延迟测量数据（最近RTT、最小RTT）用于带宽延迟积估算和拥塞程度判断；回调类型标识区分丢包事件和正常确认事件；拥塞控制状态（CA状态、CA事件、ECN状态）提供协议栈的完整状态信息。

\textbf{（2）多信号融合拥塞检测方法}

本文提出了一种多信号融合的拥塞检测方法，采用分层优先级机制综合分析多种拥塞信号。该方法的核心思想是：不同类型的拥塞信号具有不同的可靠性和紧迫程度，应当采用差异化的处理策略。具体而言，显式丢包信号具有最高优先级，表明网络已经发生严重拥塞；当拥塞控制状态机处于CA\_LOSS状态时进一步区分超时丢包事件；ECN拥塞经历（Congestion Experienced，CE）标记和ECN回显（ECN-Echo，ECE）接收作为次级信号，表明网络设备已检测到队列积压。

为了避免对瞬态扰动的过度响应，本文在拥塞检测中对不同信号类型采用差异化的触发逻辑，确保拥塞响应的精准性和鲁棒性。这种设计在保证对真实拥塞快速响应的同时，避免了因瞬态信号导致的不必要窗口缩减，有效平衡了网络稳定性和吞吐量。

\textbf{（3）差异化拥塞响应机制}

本文设计了差异化的拥塞响应机制，针对不同类型的拥塞信号采用不同的窗口保留因子（Window Retention Factor）。具体而言，丢包触发的拥塞响应采用0.70的窗口保留因子，这是最激进的响应策略，因为丢包是最严重的拥塞信号；ECN触发的拥塞响应采用0.85的窗口保留因子，保留了较大比例的窗口以维持吞吐量，因为ECN是早期拥塞信号，网络可能只是轻度拥塞；超时触发的拥塞响应采用0.50的窗口保留因子，这是最严厉的响应策略，因为超时通常意味着连续多个数据包丢失或网络路径严重中断。

此外，本文引入了连续缩减保护机制（Consecutive Decrease Protection）：当连续拥塞响应次数超过3次时，保持当前窗口不再缩减，避免因连续拥塞事件导致窗口过度收缩。这种差异化响应机制的设计理念是：拥塞信号的性质和严重程度不同，响应策略也应当有所区别。传统算法对所有拥塞信号采用统一的响应策略（如窗口减半），无法区分轻度拥塞和严重拥塞，导致在轻度拥塞时响应过度、在严重拥塞时响应不足。Lark的差异化响应机制能够更精准地匹配拥塞程度，在保证网络稳定性的同时最大化吞吐量。

\textbf{（4）强化学习驱动的自适应参数调优}

本文实现了强化学习驱动的自适应参数调优机制，根据多种因素动态调整乘性增加因子（Multiplicative Increase Factor，$\alpha$）。调整因素包括：RTT膨胀比（当前RTT与最小RTT的比值）反映网络排队延迟程度，采用三级梯度响应——当RTT膨胀比小于1.3时增大$\alpha$（步长$+0.03$），介于1.3至2.0之间时温和增大$\alpha$（步长$+0.01$），大于3.0时减小$\alpha$（步长$-0.02$）；奖励信号的指数移动平均值（EMA）反映整体控制效果的趋势，正向趋势增大$\alpha$，负向趋势减小$\alpha$；连续增长次数超过5次时给予稳定增长奖励。

$\alpha$的取值被限制在$[1.05, 1.50]$范围内，基准初始值为1.20。相比固定参数设计，这种自适应机制能够使Lark在不同网络环境下自动调整控制策略的激进程度：在低拥塞环境中快速提升$\alpha$以充分利用带宽，在高拥塞环境中降低$\alpha$以保持网络稳定。

\textbf{（5）融合控制策略}

本文采用融合控制策略计算目标拥塞窗口，将基于带宽延迟积（Bandwidth-Delay Product，BDP）估计的速率控制与基于当前窗口的保守控制相结合。BDP通过窗口化最大值过滤（Windowed Max-filter）方法估计瓶颈带宽（取最近20个投递率样本的最大值），乘以最小RTT得到。在慢启动阶段，目标窗口设为BDP的2倍，采用指数增长策略以快速探测可用带宽；在拥塞避免阶段，目标窗口按公式$V_t = \max(\alpha \times BDP, W) + \gamma \times MSS$计算，其中$\gamma = 2.0$为加性增加因子，每个ACK事件增加$\gamma$个MSS的窗口增量。

融合控制策略还引入了管道利用率（Pipeline Utilization）感知机制：当管道利用率低于70\%时，额外增加1个MSS的窗口增量以提高带宽利用效率；当管道利用率低于40\%时，进一步增加2个MSS的窗口增量以快速恢复带宽利用。这种设计确保了Lark算法能够快速适应网络带宽变化，在带宽增加时迅速提升吞吐量，在带宽减少时及时收敛。

\textbf{（6）每流独立状态管理}

本文设计了每流独立的状态管理机制，为每个TCP连接维护独立的状态数据结构。状态数据包括：带宽估计数据（窗口化最大值过滤的投递率样本队列，容量20个样本点）、RTT跟踪数据（最近RTT样本队列，容量40个样本点）、BDP估计值、自适应参数$\alpha$、拥塞计数器（连续缩减次数、连续增长次数、丢包计数、ECN计数）、吞吐量EMA等。

每流独立状态管理基于哈希表实现$O(1)$时间复杂度的状态查找与更新。其意义在于：在多流并发场景下，不同流可能经历不同的网络路径和拥塞状况，采用全局统一的状态管理会导致流间干扰；而每流独立的状态管理能够实现差异化控制，提高多流场景下的整体性能和公平性。

\textbf{（7）多线程并行化仿真架构}

本文采用多线程并行化仿真架构加速算法训练和验证。传统的ns-3网络仿真器采用单线程串行架构，在大规模网络拓扑和长时间仿真场景下效率较低。本文采用MTP-Interface（Multi-Thread Parallel Interface）多线程接口，将网络拓扑划分为多个逻辑分区，在多个线程上并行执行各分区的仿真事件，通过保守同步协议协调线程间的事件同步。实验表明，在8核CPU上可获得3-5倍的仿真加速。

\section{论文组织结构}

本文的其余部分组织如下：

第2章介绍相关工作，包括传统拥塞控制算法、基于机器学习的拥塞控制方法、显式拥塞通知机制、强化学习在网络领域的应用等。

第3章详细描述Lark算法的设计与实现，包括系统架构、观测空间设计、多信号融合拥塞检测、差异化拥塞响应、自适应参数调优、融合控制策略等核心组件。

第4章介绍实验评估，包括实验环境配置、评测指标、对比算法、实验场景设计以及详细的实验结果分析。

第5章讨论Lark算法的局限性和未来研究方向。

第6章总结全文并展望未来工作。
